\chapter{数据、备份与恢复}
\label{ch:backup}
\chapterlead{备份的价值只在恢复时体现。本章用 restic 建立一套自动备份，重点是恢复演练：没有演练过的备份，等于没有备份。}

\section{起步：第一次备份与第一次恢复}

\begin{nixcode}
environment.systemPackages = with pkgs; [ restic ];
\end{nixcode}

\begin{shellsession}
$ export RESTIC_REPOSITORY=/mnt/backup/restic
$ export RESTIC_PASSWORD_FILE=/run/secrets/restic-pass
$ restic init
$ restic backup /home/weiwen /var/lib/api
$ restic snapshots
\end{shellsession}

\texttt{init} 建立仓库，\texttt{backup} 生成快照，\texttt{snapshots} 列出已有快照。restic 按内容分块去重，因此第二次备份只存储变化部分。

恢复同样简单，但必须演练：

\begin{shellsession}
$ restic restore latest --target /tmp/restore
$ restic restore <快照ID> --target /tmp/restore --include /home/weiwen/lab
$ restic mount /mnt/tmp      # 以文件系统方式浏览历史
\end{shellsession}

\texttt{mount} 把仓库挂成一个只读目录树，可以像浏览文件一样找到某个时间点的某个文件，比逐次 \texttt{restore} 直观。

\section{三种丢失与三份副本}

需要防的丢失有三类：误删（文件被人删掉）、损坏（磁盘或文件系统损坏）、灾难（整机或机房丢失）。三类需要不同的副本位置。

\begin{booktable}{\textwidth}{J{0.769} J{1.231}}
\toprule
副本 & 防住什么 \\
\midrule
\endhead
本机快照 & 误删、改坏，恢复最快 \\
异地磁盘或 NAS & 磁盘损坏、整机故障 \\
远端对象存储 & 机房级灾难 \\
\bottomrule
\end{booktable}

“三份副本、两种介质、一份异地”是常见的最低标准。个人使用时，本机加一块异地磁盘已经能覆盖绝大多数情况。

\section{自动化}

\begin{nixcode}
services.restic.backups.daily = {
  repository = "/mnt/backup/restic";
  passwordFile = "/run/secrets/restic-pass";
  paths = [
    "/home/weiwen"
    "/var/lib/api"
  ];
  exclude = [
    "/home/weiwen/.cache"
    "/home/weiwen/lab/tmp"
  ];
  timerConfig = {
    OnCalendar = "daily";
    Persistent = true;
  };
  pruneOpts = [
    "--keep-daily 7"
    "--keep-weekly 4"
    "--keep-monthly 6"
  ];
};
\end{nixcode}

\texttt{Persistent = true} 让错过的任务在下次开机时补跑。\texttt{pruneOpts} 控制保留策略：日备留七天、周备留四周、月备留六个月，是家用与服务器的通用起点。

密码文件放在 nixpkgs 的 secrets 机制之外，由部署流程写入，不进入 store。这一点很关键：备份仓库的密码若与备份放在同一台机器上，磁盘损坏时两者一起丢。

\section{验证备份是否可用}

备份任务的退出码为 0 不代表能恢复。三项检查按成本递增。

\begin{shellsession}
$ restic check                       # 元数据一致性
$ restic check --read-data-subset=5% # 抽样读取数据块
$ restic restore latest --target /tmp/verify \
    && diff -r /tmp/verify/home/weiwen/lab ~/lab
\end{shellsession}

第一项快，适合每次备份后自动执行。第二项慢，适合每周一次。第三项是真正的恢复演练，每季度做一次，把恢复过程写成步骤清单。

\begin{warnbox}[备份的三个常见失败]
一是路径写错，备份了一个空目录，任务一直是成功状态。二是排除规则过宽，把数据目录本身排掉了。三是密码丢失，仓库无法打开。前两项靠 \texttt{restic snapshots} 与 \texttt{restic ls} 抽查即可发现，第三项要求密码至少存放在两个独立位置。
\end{warnbox}

\section{数据库的备份}

直接复制数据库文件通常得到不一致的状态。正确做法是让数据库导出一份逻辑备份，再备份这个文件。

\begin{nixcode}
services.postgresqlBackup = {
  enable = true;
  databases = [ "api" ];
  startAt = "*-*-* 03:15:00";
};
\end{nixcode}

导出的文件放在 \texttt{/var/backup/\allowbreak{}postgresql}，随后由 restic 一并备份。恢复时先建库，再用 \texttt{psql} 导入。

\section{练习}

\begin{enumerate}
  \item 初始化一个 restic 仓库，备份家目录，用 \texttt{snapshots} 确认。
  \item 删除一个文件，用 \texttt{restic mount} 找到它并恢复，记录耗时。
  \item 配置每日自动备份与保留策略，重建后用 \texttt{systemctl list-timers} 确认。
  \item 运行 \texttt{restic check\allowbreak\ -{}-read-data-subset=5\%}，记录输出与耗时。
  \item 写一份恢复步骤清单，要求他人照做能在三十分钟内恢复一份文件。
\end{enumerate}

\begin{checkpoint}
本章结束时，应能：
\begin{itemize}
  \item 用 restic 完成备份、查看、恢复与挂载浏览；
  \item 说明三类丢失与三份副本的对应关系；
  \item 配置带保留策略的自动备份；
  \item 用三项检查验证备份可恢复，并解释各自成本；
  \item 正确处理数据库这类不能直接复制的数据。
\end{itemize}
\end{checkpoint}
